\section{Bölüm sonu çözümlü problemler}

\begin{enumerate}
  \item \textbf{Eksiklik oranı.} 240 öğrencinin 18'inde başlangıç puanı,
  36'sında son-test puanı eksiktir. Değişken bazındaki eksiklik oranlarını
  bulun.

  \textbf{Çözüm:} Başlangıç için $18/240=0{,}075$, yani yüzde 7,5; son-test
  için $36/240=0{,}15$, yani yüzde 15'tir. Aynı öğrencilerin iki değişkende de
  eksik olup olmadığı bilinmeden tam gözlem sayısı yalnız bu oranlardan
  bulunamaz.

  \item \textbf{Mekanizma yorumu.} Son-test eksikliği yalnız gözlenen
  başlangıç puanı ve program grubuyla ilişkili görünmektedir. Hangi mekanizma
  çalışma varsayımı olarak düşünülebilir?

  \textbf{Çözüm:} Başlangıç puanı ve grup atama/analiz modeline alındığında MAR
  makul bir çalışma varsayımı olabilir. Gözlenmeyen son-test değerinin ayrıca
  eksikliği etkileyip etkilemediği gözlenen veriden kanıtlanamaz; MNAR
  duyarlılık analizi gerekir.

  \item \textbf{Ortalama ile atama.} Gözlenen değerler $40,50,60$ ve bir eksik
  hücre olsun. Eksik hücreye gözlenen ortalama yazıldığında ortalama ve
  örneklem varyansı nasıl değişir?

  \textbf{Çözüm:} Gözlenen ortalama 50'dir. $40,50,60,50$ dizisinin ortalaması
  yine 50 olur. İlk üç değerin örneklem varyansı 100, atama sonrasındaki
  örneklem varyansı $200/3\approx66{,}67$ olur. Yapay merkez değeri değişkenliği
  küçültür.

  \item \textbf{Rubin toplam varyansı.} $m=5$, atama içi ortalama varyans
  $\bar U=4$ ve atamalar arası varyans $B=1{,}5$ olsun. Toplam varyansı ve
  yaklaşık standart hatayı bulun.

  \textbf{Çözüm:}
  \[
  T=4+\left(1+\frac15\right)(1{,}5)=5{,}8,
  \qquad \SE=\sqrt{5{,}8}\approx2{,}41.
  \]
  Yalnız $\sqrt{\bar U}=2$ kullanmak atama belirsizliğini dışarıda bırakır.

  \item \textbf{Duyarlılık kararı.} MAR çoklu atama analizinde grup farkı 5,2
  puan; eksik sonuçların her grupta 4 puan daha düşük varsayıldığı senaryoda
  4,9 puan; yalnız müdahale grubunda 8 puan daha düşük varsayıldığı senaryoda
  1,1 puandır. Sonucu değerlendirin.

  \textbf{Çözüm:} Benzer düşüş varsayımında sonuç kararlıdır. Müdahale
  grubundaki güçlü seçici kayıp varsayımında fark büyük ölçüde azalmaktadır.
  Ana bulgu eksikliğin gruplar arasında farklı davranmadığı varsayımına
  duyarlıdır; bu sınır açıkça raporlanmalıdır.
\end{enumerate}